%%%%%%%%%%%%%%%%%%%%%%acronym.tex%%%%%%%%%%%%%%%%%%%%%%%%%%%%%%%%%%%%%%%%%
% sample list of acronyms
%
% Use this file as a template for your own input.
%
%%%%%%%%%%%%%%%%%%%%%%%% Springer %%%%%%%%%%%%%%%%%%%%%%%%%%

\extrachap{Siglas}

A continuación se presenta el listado de las siglas y acrónimos empleados a lo largo del presente documento, así como la definición formal de los conceptos subyacentes.

\begin{description}[NAND]
\item[AC]{Autómata Celular. Sistema dinámico y discreto en el espacio, el tiempo y el estado, cuya evolución es regida por reglas de transición locales.}
\item[AND]{Compuerta lógica de conjunción, definida por retornar un estado afirmativo si y solo si la totalidad de sus operandos son afirmativos.}
\item[CPU]{\textit{Central Processing Unit}. Unidad Central de Procesamiento encargada de la ejecución secuencial de las instrucciones lógicas y aritméticas en la arquitectura subyacente.}
\item[ECS]{\textit{Entity Component System}. Patrón arquitectónico de diseño de software que estructura el sistema mediante la desvinculación estricta de datos y operaciones, optimizando la iteración sobre la memoria caché.}
\item[GE]{\textit{Garden of Eden}. Configuración del espacio celular o estado de tipo Jardín del Edén, caracterizada axiomáticamente por la inexistencia de una preimagen en la función de transición global.}
\item[ID]{Identificador. Cadena alfanumérica o numérica unívoca asignada para distinguir de manera rigurosa una entidad, componente o requerimiento.}
\item[JSON]{\textit{JavaScript Object Notation}. Formato estándar y ligero de intercambio de información, empleado para la serialización y persistencia de las configuraciones de la malla discreta.}
\item[MB]{Megabyte. Unidad de medida estándar para el almacenamiento de información en sistemas digitales, equivalente a $2^{20}$ bytes.}
\item[MT]{Máquina de Turing. Modelo matemático abstracto de computación, propuesto como marco axiomático para delimitar la resolubilidad y complejidad algorítmica de los lenguajes formales.}
\item[NAND]{\textit{Not AND}. Compuerta lógica que opera como la negación de una conjunción, demostrada formalmente como un operador computacional universal.}
\item[NOT]{Compuerta lógica de negación, la cual invierte estrictamente el valor de verdad de una proposición dada.}
\item[NP]{\textit{Nondeterministic Polynomial time}. Clase de complejidad computacional conformada por el conjunto de problemas de decisión cuyas soluciones pueden ser verificadas por una máquina determinista en tiempo polinomial.}
\item[OR]{Compuerta lógica de disyunción inclusiva, definida por arrojar un valor afirmativo si al menos uno de sus operandos es afirmativo.}
\item[RF]{Requerimiento Funcional. Declaración explícita de un comportamiento, interacción o procesamiento que el sistema computacional debe ejecutar intrínsecamente por diseño.}
\item[RNF]{Requerimiento No Funcional. Restricción o atributo de calidad del sistema que delimita los criterios operativos de rendimiento, sincronismo o mantenibilidad estructural.}
\item[SAT]{\textit{Boolean Satisfiability Problem}. Problema algorítmico centrado en determinar la existencia de una asignación de variables que satisfaga íntegramente una fórmula booleana dada.}
\item[SIMD]{\textit{Single Instruction, Multiple Data}. Arquitectura de procesamiento a nivel de registro que permite la aplicación concurrente de una misma operación sobre vectores de múltiples operandos.}
\item[TPS]{\textit{Ticks Per Second}. Métrica empleada para cuantificar la tasa de convergencia temporal, definiendo el número de transiciones globales ejecutadas por el autómata celular en el transcurso de un segundo físico.}
\item[UML]{\textit{Unified Modeling Language}. Lenguaje unificado de modelado de propósito general en la ingeniería de software, empleado para la especificación visual y rigurosa de la arquitectura y la dinámica del sistema.}
\end{description}